\begin{textsample}{POS Dim 1 – 2023\_c5 – Score 23.00 – 2023\_c5/enade\_2023\_tec\_gest\_amb\_cca51bbb5c.txt}  \label{ex:f1_pos_002}
CAPÍTULO 4 DISTRIBUIÇÃO DE CONCEITOS A \textbf{finalidade} deste capítulo é apresentar um \textbf{panorama} nacional sobre a distribuição do Conceito Enade 2023.
Para traçar esse \textbf{panorama}, recorre -se a tabelas e gráficos, que são acompanhados de comentários e análises que contribuem para a contextualização dos dados neles expostos.
Para \textbf{visualizar} o \textbf{panorama} do Enade 2023, em nível nacional, na primeira seção deste capítulo, os dados sobre o número de cursos participantes são organizados por grande região.
Na \textbf{segunda}, o mapeamento do conceito do exame é feito por categoria administrativa das IES e modalidade de oferta dos cursos, de maneira \textbf{cruzada} às grandes regiões.
Na última seção, a lógica de apresentação dos dados é feita por organização acadêmica e, também, por grande região de funcionamento dos cursos, no país como um todo. 4.1 \textbf{PANORAMA} NACIONAL DA DISTRIBUIÇÃO DOS CONCEITOS A lógica adotada na apresentação dos dados, nesta seção e nas duas seguintes, pressupõe o registro de duas informações de natureza \textbf{reiterativa}.
O primeiro diz respeito aos valores possíveis do Enade, que variam em níveis de 1 a 5, sendo que, à medida que esses valores aumentam, o curso é melhor avaliado, e ao fato de que aqueles cursos que tiveram somente um ou nenhum concluinte participante não obtêm o Conceito Enade, ficando sem conceito ( SC ), como ressaltado na apresentação deste relatório.
O segundo registro é que, caso haja diferença entre os cursos \textbf{tabulados} neste capítulo e no capítulo 1, ela corresponderá aos cursos que não foram avaliados, ou seja, àqueles, em princípio, sem estudantes concluintes que tivessem sido inscritos para a prova.
Na Tabela 4.1, são apresentadas a quantidade e a distribuição de cursos de Tecnologia em Gestão Ambiental participantes do Enade 2023, por faixa de conceito e grande região.
A partir dos dados da Tabela 4.1, nota -se que, dos 150 cursos participantes, 52 ( 34,7 \% ) classificaram -se com conceito 3, o valor modal.
Esse foi também o conceito modal em duas regiões : Sudeste ( 43,5 \% ), Sul ( 37,0 \% ).
O conceito 2 foi o segundo mais frequente em nível nacional ( 23,3 \%, correspondendo a 35 cursos ), tendo sido o conceito mais frequente no Nordeste ( 32,4 \% ) e no Norte ( 58,3 \% ).
O conceito 4 foi o terceiro ( 17,3 \%, correspondendo a 26 cursos ) e foi o conceito modal 1 região : Centro-Oeste ( 33,3 \% ).
Houve, ainda, 9 ( 6,0 \% ) cursos que receberam o conceito 5 e cinco ( 3,3 \% ) cursos que receberam o conceito 1.
Dos 150 cursos de Tecnologia em Gestão Ambiental, 23 ( 15,3 \% ) ficaram sem conceito ( SC ).
A região Norte participou com 12 cursos, o que corresponde a 8,0 \% do total nacional.
Desses, sete cursos receberam o conceito 2, o que \textbf{equivale} a 58,3 \% do total regional, o conceito modal na região, conforme mencionado.
A quatro cursos ( 33,3 \% ) atribuiu -se o conceito 3 ; a um curso ( 8,3 \% ), o conceito 1 ; e a nenhum curso ( 0,0 \% ) o conceito 5.
Ainda nessa região, nenhum curso ( 0,0 \% ) obteve conceito 4 e nenhum curso ( 0,0 \% ) ficou sem conceito ( SC ).
A região Nordeste participou com 34 cursos, correspondentes a 22,7 \% do total nacional.
Nessa região, 11 cursos ( 32,4 \% em termos regionais ) obtiveram conceito 2, o conceito modal para a região, como já comentado.
Os conceitos 3 e 4 foram atribuídos, respectivamente, a oito e a cinco cursos ( 23,5 \% e 14,7 \% ).
Ainda nessa região, três cursos ( 8,8 \% ) foram avaliados com o conceito 5 e um curso ( 2,9 \% ) foi avaliado com o conceito 1.
Outros seis cursos ( 17,6 \% ) ficaram sem conceito ( SC ).
Para os 46 cursos participantes da região Sudeste ( 30,7 \% do total nacional ), o conceito 3 foi o conceito modal, como já anteriormente assinalado, tendo sido obtido por 20 cursos ( 43,5 \% ).
O conceito 4 foi atribuído a oito cursos ( 17,4 \% ) e o conceito 2, a oito cursos ( 17,4 \% ).
Receberam os conceitos 5 e 1, respectivamente, três cursos ( 6,5 \% ) e um curso ( 2,2 \% ).
Nessa região, seis cursos ( 13,0 \% ) ficaram sem conceito ( SC ).
Os 46 cursos da região Sul corresponderam a 30,7 \% do total nacional.
Como já \textbf{indicado}, houve predominância do conceito 3 ( conceito modal ), atribuído a 37,0 \% dos cursos da região, ou seja, atribuído a 17 dos 46 cursos participantes na região Sul.
O conceito 4 foi atribuído a nove cursos ( 19,6 \% ) ; o conceito 2, sete cursos ( 15,2 \% ) ; e o conceito 5, a dois cursos ( 4,3 \% ).
Na região Sul, dois cursos ( 4,3 \% ) obtiveram conceito 1 e nove cursos ( 19,6 \% ) ficaram sem conceito ( SC ).
Já dos 12 cursos participantes na região Centro-Oeste ( 8,0 \% do total nacional ), 4 ( 33,3 \% em termos regionais ) receberam conceito 4, o conceito modal, como já destacado.
Três cursos ( 25,0 \% ) obtiveram o conceito 3, dois cursos ( 16,7 \% ) obtiveram o conceito 2 e um curso ( 8,3 \% ), o conceito 5.
Nessa região, nenhum curso ( 0,0 \% ) obteve o conceito 1 e dois cursos ( 16,7 \% ) ficaram sem conceito ( SC ).
Em um gráfico de distribuição \textbf{cumulativa}, em cada valor de \textbf{abscissa} ( eixo x ), o valor apresentado no eixo das \textbf{ordenadas} ( eixo y ) é o valor acumulado da \textbf{variável} ( Conceito Enade ).
Em linhas gerais, em um gráfico desta mesma \textbf{tipologia}, \textbf{poligonais} mais à esquerda representam, no caso dos dados apresentados nesta seção, as grandes regiões com a pior distribuição, e \textbf{poligonais} mais à direita, as grandes regiões com os melhores conceitos.
A partir destas diretrizes, no Gráfico 4.1, é apresentada a distribuição \textbf{cumulativa} do Conceito Enade, segundo a grande região.
No Gráfico 4.1, observando a linha vertical \textbf{pontilhada}, que ajuda a \textbf{visualizar} o percentual de conceitos acumulados até o valor 3, ou seja, conceitos 1, 2 e 3, pode -se notar, por exemplo, que a região Sudeste ( linha vermelha ) apresenta 72,5 \% dos cursos com conceito 3 ou abaixo deste valor.
A região Norte ( linha azul ), por outro lado, apresenta 100,0 \% dos cursos na mesma situação : com conceito 3 ou abaixo.
Destaca -se que, a região Norte, os conceitos 4 e 5 não foram alcançados por nenhum curso, na qual todos os cursos têm conceito menor ou igual a 3, e a respectiva \textbf{poligonal} \textbf{cumulativa} alcançou 100 \% no conceito 3.
A região Norte apresenta os cursos com conceitos mais baixos, ao passo que a região Centro-Oeste apresenta a distribuição com valores maiores.
Os cursos em IES situadas nas regiões Sul ( linha \textbf{roxa} ), Nordeste ( linha verde ) e Sudeste ( linha vermelha ) apresentam uma situação intermediária. 4.2 CONCEITOS POR CATEGORIA ADMINISTRATIVA, MODALIDADE DE OFERTA DOS CURSOS E GRANDE REGIÃO Na Tabela 4.2, é apresentada a distribuição dos cursos de Tecnologia em Gestão Ambiental participantes do Enade 2023, por categoria administrativa e modalidade de oferta, de acordo com os conceitos por eles alcançados, segundo as grandes regiões brasileiras.
De acordo com as informações da Tabela 4.2, dos 150 cursos participantes, 58 ( 38,7 \% ) eram \textbf{ministrados} em instituições públicas e 92 ( 61,3 \% ), em instituições privadas.
Quanto à modalidade de oferta, oferta, 67 ( 44,7 \% ) cursos eram presenciais e 83 ( 55,3 \% ), a distância.
De acordo com as informações da mesma tabela, em termos nacionais, entre instituições públicas e privadas, observa -se que, dos 9 cursos avaliados com conceito 5, 4 foram oferecidos em IES públicas, e 5, em IES privadas.
Dos 92 cursos participantes de IES privadas, o conceito 3 foi o valor modal, atribuído a 31 ( 33,7 \% ) cursos.
Em relação aos demais cursos de instituições privadas, 18 cursos ( 19,6 \% ) foram avaliados com o conceito 2, 13 cursos ( 14,1 \% ) foram avaliados com o conceito 4, cinco cursos ( 5,4 \% ) foram avaliados com o conceito 5 e, três cursos ( 3,3 \% ), com o conceito 1.
Nessa categoria, 22 cursos ( 23,9 \% ) ficaram sem conceito ( SC ).
Na rede pública, o conceito modal foi o 3, atribuído a 21 cursos, correspondentes a 36,2 \% dos 58 cursos da categoria.
Dos demais cursos participantes, 17 cursos ( 29,3 \% ) obtiveram o conceito 2, 13 cursos ( 22,4 \% ) receberam o conceito 4 e quatro cursos ( 6,9 \% ), o conceito 5.
O conceito 1 foi atribuído a dois cursos ( 3,4 \% ).
Nessa categoria administrativa, um curso ( 1,7 \% ) ficou sem conceito ( SC ).
Dos 67 cursos participantes da modalidade presencial, o conceito 3 foi o valor modal, atribuído a 22 ( 32,8 \% ) cursos.
Dos demais cursos dessa modalidade de oferta, 17 cursos ( 25,4 \% ) receberam o conceito 2 ; 14 cursos ( 20,9 \% ) receberam o conceito 4 ; quatro cursos ( 6,0 \% ), o conceito 5 ; e três cursos ( 4,5 \% ), o conceito 1.
Nessa modalidade de oferta, sete cursos ( 10,4 \% ) ficaram sem conceito ( SC ).
Dos 83 cursos participantes da modalidade a distância, o conceito 3 foi o valor modal, atribuído a 30 ( 36,1 \% ) cursos.
Dos demais cursos \textbf{ofertados} nessa modalidade, 18 cursos ( 21,7 \% ) receberam o conceito 2, 12 cursos ( 14,5 \% ) receberam o conceito 4 ; cinco cursos ( 6,0 \% ) receberam o conceito 5 ; e dois cursos ( 2,4 \% ) receberam conceito 1.
Nessa modalidade de oferta, 16 cursos ( 19,3 \% ) ficaram sem conceito ( SC ).
Na análise por região, observa -se, que, na região Norte, as instituições públicas participaram com sete cursos ( 58,3 \% total regional ), dos quais nenhum ficou sem conceito ( SC ) ou recebeu os conceitos 4 e 5.
Os conceitos 2 e 3 foram atribuídos a três cursos cada, correspondendo cada um deles a 42,9 \% dessa categoria na região.
O conceito 1 teve um curso ( 14,3 \% ).
As instituições privadas participaram com 5 cursos na região Norte ( 41,7 \% do total regional ), e nenhum curso ficou sem conceito ( SC ).
O conceito modal nas instituições privadas na região foi o 2, atribuído a quatro cursos, o que corresponde a 80,0 \% dessa categoria na região.
O conceito 3 foi atribuído a um curso.
Nenhum curso ficou sem conceito ( SC ) ou recebeu os demais conceitos.
Dos cursos oferecidos na região Norte, sete eram presenciais e nenhum ficou sem conceito ( SC ).
Dos sete cursos oferecidos, três obtiveram o conceito 3 ( modal ) ; três, o conceito 2 ; um, o conceito 1.
Nenhum curso ficou sem conceito ( SC ) ou recebeu os demais conceitos na região Norte e na modalidade presencial.
Cinco cursos eram na modalidade a distância, tendo quatro recebido o conceito modal 2 ; um, o conceito 3.
Nenhum curso ficou sem conceito ( SC ) ou recebeu os demais conceitos na região Norte e na modalidade a distância.
Na região Nordeste, a rede privada concentra 17 dos 34 cursos participantes, o equivalente a 50,0 \% do total da região.
Entre os cursos oferecidos por IES privadas no Nordeste, quatro cursos foram avaliados com conceito 3, o conceito modal ; três cursos obtiveram o conceito 2 ; dois cursos, o conceito 5 ; um curso, o conceito 4 ; e um curso obteve o conceito 1.
Nessa categoria administrativa e grande região, seis cursos ficaram sem conceito ( SC ).
As instituições públicas da região Nordeste participaram com 17 cursos ( 50,0 \% do total da região ).
Desses, oito ( 47,1 \% ) obtiveram conceito 2, o valor modal, seguidos por quatro cursos, que receberam o conceito 4.
O conceito 3 foi atribuído a quatro cursos ; o conceito 5, a um curso a nenhum curso.
Nenhum curso ficou sem conceito ( SC ) ou recebeu o conceito 1.
Dos cursos oferecidos na região Nordeste, todos ( 21 ) eram presenciais : três ficaram sem conceito ( SC ), um obteve o conceito 1 ; oito, o conceito 2 ; quatro, o conceito 3, o valor modal ; quatro, o conceito 4 ; e um, o conceito 5.
Um total de 13 cursos eram \textbf{ofertados} na modalidade a distância, tendo quatro cursos recebido o conceito modal 3 ; um, o conceito 4 e dois, o conceito 5.
Três cursos ficaram sem conceito ( SC ).
Nenhum curso recebeu o conceito 1.
Na região Sudeste, a proporção de cursos da rede privada, 63,0 \%, correspondeu a 29 dos 46 cursos participantes.
Nessa combinação de categoria administrativa e grande região, o conceito modal foi o 3, atribuído a 12 cursos ; e cinco cursos ficaram sem conceito ( SC ).
Os demais foram avaliados com os seguintes conceitos : 2 ( seis cursos ), 4 ( quatro cursos ), 5 ( um curso ) e 1 ( um curso ).
Dos 17 cursos oferecidos por instituições públicas na região Sudeste, dois obtiveram o conceito 5 ; quatro, o conceito 4 ; oito, o conceito 3 ; dois obtiveram conceito 2 e nenhum obteve o conceito 1.
Um curso ficou sem conceito ( SC ).
Ainda na região Sudeste, menos da metade dos cursos ( 43,5 \% ) foi oferecida na modalidade presencial : 20 dos 46 cursos.
Para esses, a moda foi o conceito 3, com oito cursos.
Os demais cursos presenciais da região foram alocados nos conceitos 4 ( cinco cursos ), 5 ( dois cursos ), 2 ( dois cursos ) e 1 ( nenhum curso ).
Três cursos presenciais ficaram sem conceito ( SC ).
Na região Sudeste, 26 cursos foram oferecidos na modalidade a distância, tendo sido um com conceito 1 ; seis com conceito 2 ; 12 com conceito 3 ; três com conceito 4 ; e um com conceito 5.
Três cursos ficaram sem conceito ( SC ).
As instituições privadas concentraram 33 dos 46 cursos participantes da região Sul, 71,7 \% do total regional.
Desses, 12 cursos obtiveram conceito 3, o conceito modal.
Nessa combinação de categoria administrativa e grande região, cinco receberam o conceito 4 ; cinco receberam o conceito 2 ; um, o conceito 5 ; e um obteve o conceito 1.
Nove cursos ficaram sem conceito ( SC ).
As instituições públicas na região Sul participaram com 13 cursos ( 28,3 \% dos cursos da região ), aos quais foram atribuídos os conceitos 1 ( um curso ), 2 ( dois cursos ), 3 ( cinco cursos ), 4 ( quatro cursos ) e 5 ( um curso ).
Nenhum curso da região Sul oferecido por IES públicas ficou sem conceito ( SC ).
Na região Sul, 14 dos cursos foram oferecidos na modalidade presencial.
Para esses, a moda foi o conceito 3, com 5 cursos.
Os demais cursos presenciais da região foram alocados nos conceitos 4 ( quatro cursos ), 2 ( dois cursos ), 5 ( um curso ) e 1 ( um curso ).
Ficou sem conceito ( SC ) um curso.
Na modalidade a distância, foram oferecidos 32 cursos, sendo um curso com conceito 1, cinco cursos com conceito 2, 12 cursos com conceito 3 ( modal ), cinco cursos com conceito 4 e um curso com conceito 5.
Oito cursos ficaram sem conceito ( SC ) nessa região.
Na região Centro-Oeste, 8 dos 12 cursos participantes eram de IES privadas ( 66,7 \% em termos regionais ).
Desses, três cursos receberam conceito 4, o valor modal.
Dos demais cursos, dois receberam o conceito 3, um recebeu o conceito 5, nenhum curso recebeu os conceitos 1 e 2.
Dois cursos ficaram sem conceito ( SC ).
Quanto aos quatro cursos oferecidos por instituições públicas na região Centro-Oeste ( 33,3 \% do total regional ), dois ficaram com conceito 2 ; um, com conceito 3 ; um, com conceito 4 ; e nenhum ficou com conceitos 1 e 5 ou sem conceito ( SC ) nessa região.
Também na região Centro-Oeste, na modalidade presencial, foram \textbf{ofertados} 5 ( 41,7 \% ).
Para esses, a moda foi o conceito 2, com 2 cursos ( 40,0 \% ).
Os demais cursos presenciais da região foram alocados nos conceitos 2 ( dois cursos, 40,0 \% ), 4 ( um curso, 20,0 \% ), 5 ( nenhum curso, 0,0 \% ) e 1 ( nenhum curso, 0,0 \% ).
Nenhum curso presencial da região ficou sem conceito ( SC ).
A maioria dos cursos foram \textbf{ofertados} na modalidade a distância, com sete cursos, sendo um curso com conceito 3 ( modal ), três cursos com conceito 4 e um curso com conceito 5.
Dois cursos ficaram sem conceito ( SC ) nessa região.
Nenhum curso recebeu os conceitos 1 e 2.
A título de síntese, no Gráfico 4.2, apresenta -se a distribuição do Conceito Enade, segundo a categoria administrativa da IES.
Nesse gráfico, como dito anteriormente, \textbf{poligonais} mais à esquerda representam Conceito Enade mais baixos, e \textbf{poligonais} mais à direita, Conceito Enade mais altos.
É importante observar que, para cada \textbf{poligonal} presente nesse gráfico, a porcentagem constante no eixo y refere -se ao universo de cursos oferecidos por instituições públicas ( \textbf{poligonal} azul ) ou privadas ( \textbf{poligonal} verde ), respectivamente.
Os cursos em IES públicas ( \textbf{poligonal} azul ) apresentam uma \textbf{poligonal} mais à direita do que os em IES privadas ( \textbf{poligonal} verde ) e, consequentemente, uma distribuição de Conceito Enade com valores maiores. 4.3 CONCEITOS POR ORGANIZAÇÃO ACADÊMICA E GRANDE REGIÃO Na Tabela 4.3, encontra -se a distribuição dos conceitos atribuídos aos cursos participantes do Enade 2023, na área de Tecnologia em Gestão Ambiental, por organização acadêmica, segundo as grandes regiões brasileiras.
Como mostram os dados da Tabela 4.3, dos 150 cursos de Tecnologia em Gestão Ambiental participantes, 46 eram oferecidos em universidades, 51 em centros universitários e 15 em faculdades ; 38 cursos foram oferecidos por CEFET/IF.
Essa distribuição corresponde a, respectivamente, 30,7 \%, 34,0 \%, 10,0 \% e 25,3 \% dos cursos.
De acordo com os dados apresentados na mesma tabela, de todos os 26 cursos avaliados com o conceito 4, 9 eram vinculados a universidades.
Os cursos vinculados a esse tipo de organização acadêmica tiveram o conceito 3 como conceito modal, atribuído a 15 cursos ( 32,6 \% ).
Os demais cursos \textbf{ofertados} por universidades receberam os conceitos 2 ( nove cursos ), 5 ( quatro cursos ) e 1 ( dois cursos ).
Sete cursos ficaram sem conceito ( SC ).
Dos cursos vinculados a centros universitários, o conceito modal também foi o 3, atribuído a 17 cursos ( 33,3 \% ).
Treze cursos vinculados a esse tipo de organização acadêmica ficaram sem conceito ( SC ), e os outros cursos receberam os conceitos 2 ( 11 cursos ), 4 ( seis cursos ), 5 ( dois cursos ) e 1 ( dois cursos ).
Dos 15 cursos mantidos por faculdades, 4 ( 26,7 \% ) receberam conceito 3, o conceito modal.
O segundo conceito mais frequentemente atribuído aos cursos oferecidos por essa categoria de IES foi o conceito 3, recebido por quatro cursos.
Na sequência, três cursos obtiveram conceito 4 ; um curso, o conceito 5 ; e nenhum curso, o conceito 1.
Dos cursos oferecidos em faculdades, três cursos ficaram sem conceito ( SC ).
Dos cursos oferecidos por CEFET/IF, o conceito modal foi o 3, atribuído a 16 cursos ( 42,1 \% ).
Nenhum curso vinculado a esse tipo de organização acadêmica ficou sem conceito ( SC ), e os outros cursos receberam os conceitos 2 ( 11 cursos ), 4 ( oito cursos ), 5 ( dois cursos ) e 1 ( um curso ).
Considerando -se, separadamente, as regiões brasileiras, verifica -se que, na região Norte, as universidades não tiveram nenhum curso avaliado.
Os centros universitários da região Norte foram representados por quatro cursos, dos quais três cursos obtiveram o conceito 2, o valor modal ; um curso, o conceito 3.
Nenhum curso obteve os conceitos 1, 4 e 5 ou ficou sem conceito ( SC ).
As faculdades participaram com um curso na região Norte, o qual recebeu o conceito 2.
Os CEFET/IF participaram com sete cursos, aos quais foram atribuídos os seguintes conceitos : 1 a um curso ; 2 a três cursos ; 3 a três cursos.
Nenhum curso obteve os conceitos 4 e 5 ou ficou sem conceito ( SC ).
Na região Nordeste, as universidades participaram com 8 dos 34 cursos da área de Tecnologia em Gestão Ambiental oferecidos.
O conceito modal foi 3, atribuído a três cursos.
Outros três cursos obtiveram o conceito 2 ; um curso, o conceito 5.
Um curso ficou sem conceito ( SC ).
Nenhum curso obteve os conceitos 1 e 4.
Os centros universitários \textbf{contaram} com dez cursos participantes na região Nordeste, três dos quais receberam o conceito modal 3.
Para os demais cursos, foram atribuídos os conceitos 5 ( um curso ), 4 ( um curso ), 2 ( um curso ) e 1 ( um curso ).
Três cursos oferecidos por centros universitários no Nordeste ficaram sem conceito ( SC ).
As faculdades foram representadas por três cursos na região Nordeste, os quais receberam os conceitos 2 ( um curso ).
Dois cursos ficaram sem conceito ( SC ).
Nenhum curso obteve os conceitos 1, 3, 4 e 5.
Na região Nordeste, os CEFET/IF participaram com 13 cursos, aos quais foram atribuídos os seguintes conceitos : 2 a seis cursos ; 3 a dois cursos ; 4 a quatro cursos ; e 5 a um curso.
Nenhum curso obteve o conceito 1 e, nenhum ficou sem conceito ( SC ).
Na região Sudeste, as universidades concentraram 16 dos 46 cursos de Tecnologia em Gestão Ambiental da região.
Dos cursos oferecidos em IES com este tipo de organização acadêmica, nessa região, o conceito modal foi 3, atribuído a seis cursos.
Os demais cursos receberam os conceitos 2 ( cinco cursos ) e 4 ( três cursos ).
Dois cursos ficaram sem conceito ( SC ).
Nenhum curso obteve os conceitos 1 e 5.
Os centros universitários participaram com 13 cursos na região Sudeste, dos quais seis obtiveram o conceito modal 3.
O conceito 5 foi atribuído a um curso ; o conceito 4 a um curso ; o conceito 2 a um curso ; e o conceito 1 a um curso.
Três cursos ficaram sem conceito ( SC ) nessa combinação de organização acadêmica e região.
Já as faculdades foram representadas por oito cursos na região Sudeste, dos quais um ficou sem conceito ( SC ).
O conceito modal 3 foi atribuído a 3 cursos.
Aos demais cursos foram atribuídos os conceitos 2 ( dois cursos ), 5 ( um curso ), 4 ( um curso ) e 1 ( nenhum curso ).
Os CEFET/IF participaram com nove cursos na região Sudeste, aos quais foram atribuídos os seguintes conceitos : 3 a cinco cursos ; 4 a três cursos ; e 5 a um curso.
Nenhum curso obteve os conceitos 1 e 2 ou ficou sem conceito ( SC ).
Dos 46 cursos da região Sul, 18 eram mantidos por universidades, para os quais o conceito modal foi 3, atribuído a 6 cursos.
Os demais cursos receberam os conceitos 4 ( cinco cursos ), 5 ( dois cursos ), 1 ( dois cursos ) e 2 ( um curso ).
Em IES com esse tipo de organização acadêmica, dois cursos da região Sul ficaram sem conceito ( SC ).
Dos 19 cursos participantes de centros universitários da região Sul, sete ficaram sem conceito ( SC ), e cinco cursos receberam conceito 3.
Os conceitos \textbf{restantes} ficaram assim distribuídos : 2 ( cinco cursos ), 4 ( dois cursos ) ; e nenhum curso recebeu os conceitos 1 e 5.
Quanto aos dois cursos vinculados a faculdades na região Sul, um curso recebeu o conceito 3 e o outro o conceito 4.
Os CEFET/IF participaram com sete cursos na região Sul, aos quais foram atribuídos os seguintes conceitos : 2 a um curso ; 3 a cinco cursos ; 4 a um curso.
Nenhum obteve os conceitos 1 e 5 ou ficou sem conceito ( SC ).
Na região Centro-Oeste, 4 dos 12 cursos oferecidos eram mantidos por universidades.
Para esse tipo de organização acadêmica, o conceito modal foi 5, atribuído a um curso.
Em relação às universidades, dois cursos da região ficaram sem conceito ( SC ) e um curso recebeu o conceito 4.
Os centros universitários da região Centro-Oeste \textbf{contaram} com cinco cursos : dois cursos foram avaliados com o conceito 4 ( modal ) ; dois cursos, com o conceito 3 ; um curso, com o conceito 2 ; e nenhum curso obteve o conceito 1 ou 5 ou ficou sem conceito ( SC ).
Houve um curso oferecido em faculdades na região Centro-Oeste, que recebeu o conceito 4 ( modal ).
No que se refere aos CEFET/IF da região Centro-Oeste, dois cursos foram avaliados : um curso deles com o conceito 3 ( modal ) ; um curso com o conceito 2 ; e nenhum curso ficou sem conceito ( SC ) ou obteve os conceitos 1, 4 e 5.
De maneira semelhante à forma como se \textbf{procedeu} em relação à seção anterior, esta seção encerra -se com o Gráfico 4.3, que apresenta a distribuição do Conceito Enade dos cursos de Tecnologia em Gestão Ambiental segundo a organização acadêmica da IES.
No Gráfico 4.3, apresenta -se a distribuição do Conceito Enade segundo a organização acadêmica da IES.
A leitura desse último gráfico é semelhante à dos anteriores, seguindo -se a lógica de que \textbf{poligonais} deslocadas para a esquerda correspondem aos conceitos mais baixos ( 1, 2 e 3 ) e \textbf{poligonais} deslocadas para a direita correspondem a Conceitos Enade mais altos.
Nas universidades ( linha \textbf{rosa} ), todos os cursos obtiveram conceito 5 ou abaixo disso, já que a \textbf{poligonal} \textbf{cumulativa} alcança 100 \% para o conceito 5.
Além disso, os cursos de universidades apresentam uma \textbf{poligonal} mais à direita do que as demais, \textbf{denotando} uma melhor distribuição de Conceitos Enade.
No outro extremo, encontram -se os cursos mantidos por centros universitários ( linha verde ), apresentando as respectivas \textbf{poligonais} mais à esquerda e, consequentemente, uma distribuição de Conceitos Enade com os menores valores.
Os cursos mantidos por faculdades ( linha vermelha ) e CEFET/IF ( linha laranja ) apresentam \textbf{poligonais} mais centrais e, consequentemente, uma distribuição de Conceitos Enade com os valores intermediários.

% matched lemmas: abscissa, contar, cruzar, cumulativo, denotar, equivaler, finalidade, indicado, ministrar, ofertar, ordenado, panorama, poligonal, pontilhado, proceder, reiterativo, restante, rosa, roxa, segunda, tabular, tipologia, variável, visualizar
\end{textsample}
